\section{TensorLift Pipeline}
\label{sec:pipeline}

Figure~\ref{fig:overview} (lower half) decomposes TensorLift into three sequential stages that, taken together, transform an unannotated RTL design into a TAIDL specification consumable by the ACT ecosystem.
\textbf{Stage~1 (RTL-to-MLIR extraction, \S\ref{subsec:stage1})} uses an extended autoGenILA backend to symbolically unroll the design and emit, for every (instruction, ASV) pair, a low-level MLIR function written in the \texttt{arith}, \texttt{memref}, \texttt{scf}, and \texttt{affine} dialects.
\textbf{Stage~2 (semantic lifting, \S\ref{subsec:stage2})} runs an eight-pass MLIR rewrite pipeline that progressively normalizes bit-manipulation idioms, recovers multiply--accumulate and saturation structure, specializes control modes, reconstructs MAC chains as \texttt{scf.for} reductions, and lifts dot-product loops into the \texttt{linalg} dialect; a final metadata-emission pass annotates the lifted MLIR with structured \texttt{taidl.*} attributes.
\textbf{Stage~3 (TAIDL assembly, \S\ref{subsec:stage3})} consumes those annotations through a generic Python emitter that materializes a TAIDL specification with XLA-HLO-style semantics, together with the data-model declarations and instruction encodings.
The boundaries between stages are deliberately deep: Stage~1 owns all hardware-level reasoning and produces a behavior-equivalent yet still-bit-level model; Stage~2 owns all rewriting from bit-level to tensor-level semantics and is purely an MLIR pass pipeline; Stage~3 is a syntactic projection into the surface TAIDL language.
The remainder of this section walks through each stage in turn.

\subsection{Stage 1: RTL-to-MLIR Extraction}
\label{subsec:stage1}

TensorLift's first stage builds on autoGenILA~\cite{autogenilaiccad,autogeniladate}, a two-pass framework for recovering instruction-level behavior from synthesizable RTL.
Given a flattened Verilog design and a small per-instruction harness file that pins control inputs to a particular opcode while leaving data inputs symbolic, autoGenILA first uses path-sensitive taint analysis~\cite{autogenilaiccad} to identify the subset of registers that are \emph{architectural state variables} (ASVs)---registers whose values persist across instruction boundaries and form the programmer-visible state.
It then processes one (instruction, ASV) pair at a time~\cite{autogeniladate}: it symbolically unrolls the design over a user-supplied cycle bound, walks the resulting AST and constraint DAG with handlers for arithmetic, comparison, ITE/case mux, bit-select, concatenation, and submodule-call constructs, and emits an LLVM IR function whose arguments are the instruction inputs and current ASV values, and whose return value is the next-state value of the chosen ASV.
Together, the two passes yield a per-instruction architectural model that is bit-equivalent to the RTL by construction, no longer dependent on simulator infrastructure, and amenable to formal reasoning.

Although autoGenILA is the right starting point---it solves the \emph{extraction} problem with provable bit-equivalence---its output is unsuitable for tensor-level reasoning in two ways.
First, the LLVM IR backend is structurally tuned for compiled execution: it lowers conditional updates into branches and \texttt{phi} nodes, replaces typed memory accesses with pointer arithmetic, and---most consequentially---discards the RTL signal-name provenance that downstream lifting passes need to recognize buffers semantically (``which scalar argument carries the activation tile, and which carries the accumulator?'').
Second, autoGenILA was designed and evaluated on scalar RISC-V cores, and several construct families that arise routinely in DNN accelerators---signed division and modulo for DMA address arithmetic, and update functions whose intermediate constraint width does not match the declared return width---either trip the constraint dispatcher or produce verifier-rejected IR.
Both shortcomings have to be repaired before semantic lifting can proceed.

\paragraph{TensorLift's modifications.}
We retain autoGenILA's taint-based ASV inference, symbolic unrolling, and constraint-DAG traversal entirely; the modifications consist of a new code generator that replaces the LLVM IR backend, plus a small automation layer wrapped around the extraction. The original LLVM IR backend remains available unchanged and is reused as a verification oracle (\S\ref{sec:evaluation}). Concretely, the changes are:

\begin{enumerate}
\item \textbf{An MLIR-producing update-function backend.}
We add an alternative backend, \texttt{MLIRUpdateFunctionGen}, that mirrors autoGenILA's existing constraint-handler dispatch---two-operand and one-operand arithmetic, ITE/case/switch mux, bit-select and concatenation, dynamic memory selection, and submodule-call constructs---but emits MLIR in the \texttt{arith}, \texttt{memref}, \texttt{scf}, and \texttt{affine} dialects rather than LLVM IR.
The choice of MLIR is not cosmetic: the \texttt{scf} dialect preserves \texttt{if}/\texttt{else} regions instead of branching, \texttt{memref} preserves indexed loads and stores instead of pointer arithmetic, and the dialect-aware rewrite framework~\cite{mlir-cgo2021,mlir-pattern-rewriting} underpins every pass in Stage~2 (\S\ref{subsec:stage2}).
The backend is selected by a single CLI flag, shares all of autoGenILA's parsing, taint, and traversal infrastructure unchanged, and every emitted file is round-tripped through \texttt{mlir-opt} as a sanity check before lifting.

\item \textbf{Time-indexed input memrefs.}
The natural translation of autoGenILA's input model would lower each module input signal at each unroll cycle as an independent scalar argument, so a 36-cycle unroll of a controller with eleven inputs would produce nearly four hundred unrelated scalars.
Our backend instead packs the time series of each input into a single \texttt{memref}-typed argument whose leading dimension is the unroll bound and whose element type is the original signal width, with one explicit \texttt{memref.load} per cycle inside the function body.
This grouping carries no semantic content of its own, but it is what later allows the loop-reconstruction pass to recognize MAC chains as indexable reductions and lift them into \texttt{scf.for} loops (\S\ref{subsec:stage2}); without the grouping, the same chain decomposes into hundreds of unrelated scalars and is unrecoverable.

\item \textbf{RTL signal-name provenance via the \texttt{taidl.arg\_names} attribute.}
For each emitted \texttt{func.func}, the backend attaches the ordered list of original RTL signal names (\texttt{io\_in\_a\_0}, \texttt{io\_srams\_read\_0\_resp\_bits\_data}, $\ldots$) as an \texttt{ArrayAttr} on the function operation.
The lifting passes downstream consult this attribute to classify arguments as data buffers, control attributes, or accumulator state, and the TAIDL assembler in Stage~3 uses the same attribute to derive instruction-operand names matching the convention used by hand-written TAIDL specifications~\cite{taidl-ae-micro2025}.
Without this hook, every shred of semantic information about RTL signals would be lost the moment the constraint DAG is lowered into anonymous MLIR \texttt{Value}s.

\item \textbf{Robustness fixes for accelerator construct families.}
Two construct families absent from autoGenILA's published evaluation but ubiquitous in DNN accelerators required completing the dispatcher.
The first is signed and unsigned division and modulo, which appear in the DMA stride and pixel-repeat address arithmetic of Gemmini's load and store controllers; we add the corresponding \texttt{arith.divsi}/\texttt{remsi} (and unsigned-variant) cases to the constraint handler.
The second is constraint expressions whose inferred intermediate width exceeds the declared return width of an update function---a common situation when a width-conservative \texttt{i32} traversal ultimately writes a one-bit predicate register; we insert an explicit \texttt{trunci}/\texttt{extui} adapter at the function return so the emitted MLIR satisfies the verifier under all combinations of intermediate and declared widths.

\item \textbf{An automation harness around the extraction.}
Vanilla autoGenILA expects the user to hand-author four ancillary files per design---a Yosys flatten script, a reset script, a per-instruction harness (\texttt{instr.txt}), and an ASV target list (\texttt{allowed\_target.txt})---together with knowledge of clock and reset signal names, pipeline depths, and per-opcode control vectors.
We replace this with a small set of generators that, given only a directory of SystemVerilog files and a top-module name, automatically (a)~run Yosys with FIRRTL-aware assertion stripping (\texttt{chformal -remove}, required because Chisel \texttt{assert} calls survive into the synthesized netlist) and post-process its output to satisfy autoGenILA's parser, (b)~enumerate ASV targets by walking flip-flop declarations and grouping vector slices, and (c)~discover instruction encodings from the RTL decoder for funct-field ISAs.
The harness reduces the per-design manual input to two items, an RTL directory and a top-module name, which is what makes the end-to-end pipeline of Figure~\ref{fig:overview} viable in practice.
\end{enumerate}

The result of Stage~1 is a corpus of MLIR files, one per (instruction, ASV) pair, each containing a single \texttt{func.func} carrying the \texttt{taidl.arg\_names} attribute, written entirely in the \texttt{arith}, \texttt{memref}, \texttt{scf}, and \texttt{affine} dialects, and provably equivalent (modulo the LLVM IR equivalence proofs of \S\ref{sec:evaluation}) to autoGenILA's bit-accurate model of the RTL.
As Figure~\ref{fig:code} (middle panel) illustrates, even a single Gemmini PE produces 686 lines of MLIR at this stage---an artifact of the bit-level structure of autoGenILA's traversal, in which the core multiply--accumulate is buried under 383 sign-extension and bit-packing operations.
Recovering tensor-level structure from this representation is the job of Stage~2.

\subsection{Stage 2: Semantic Lifting}
\label{subsec:stage2}

Stage~2 is the analytical core of TensorLift: an eight-pass MLIR pipeline, registered behind a single \texttt{tensorlift-opt} tool, that turns the bit-accurate Stage~1 output into MLIR whose structure mirrors tensor-level semantics. The passes fall into four phases (2\,+\,3\,+\,2\,+\,1 passes respectively): (\textbf{A}) bit-manipulation canonicalization, (\textbf{B}) idiom detection for MAC, control-mode, and saturation patterns, (\textbf{C}) loop and tensor-shape reconstruction, and (\textbf{D}) TAIDL metadata emission. Table~\ref{tab:lifting-passes} lists the eight passes; the remainder of this subsection highlights the design choices that distinguish the pipeline from a conventional MLIR lowering.

\begin{table}[t]
\centering
\caption{TensorLift's eight-pass MLIR semantic-lifting pipeline. Each pass is named with its \texttt{tensorlift-opt} flag.}
\label{tab:lifting-passes}
\setlength{\tabcolsep}{3pt}
\renewcommand{\arraystretch}{1.18}
\scriptsize
\begin{tabular}{@{}l p{5.9cm}@{}}
\toprule
Phase / Pass & What it does \\
\midrule
A1 \texttt{canon-bitmanip}
& Collapses the bit-by-bit sign-extension chains that autoGenILA emits when traversing Verilog \texttt{\$signed} contexts into a single \texttt{arith.extsi}; the dominant source of code reduction on PEs. \\
A2 \texttt{narrow-types}
& Folds redundant \texttt{trunci}/\texttt{ext} round trips left over after canonicalization, while deliberately preserving the \texttt{extsi(trunci(\,))} idioms that Pass~B5 must recover as saturation. \\
\addlinespace[1pt]
B3 \texttt{detect-mac}
& Walks each \texttt{addi} back through width casts to find a multiplier and recovers its pre-extension inputs; tags the pair with \texttt{tensorlift.mac} when the operand widths are hardware-realistic, filtering out bit-packing artifacts. \\
B4 \texttt{specialize-control}
& Constant-folds the loads of the fixed control inputs of the instruction under analysis (taken from the same \texttt{instr.txt} that drove Stage~1), letting downstream canonicalization eliminate the \texttt{select} chains that selected between hardware modes. \\
B5 \texttt{detect-clamp}
& Recognizes the hardware fixed-point saturation idiom \texttt{ext(trunci(x))} and annotates the surviving extension with the recovered clamp range and signedness. \\
\addlinespace[1pt]
C6 \texttt{reconstruct-loops}
& Builds a use--def chain over MAC-annotated additions whose accumulator input is the previous MAC's output, and---for chains of length~$\geq\!2$---materializes the chain as an \texttt{scf.for} reduction with a single \texttt{iter\_arg}. \\
C7 \texttt{lift-to-linalg}
& Verifies that a reconstructed \texttt{scf.for} matches the canonical dot-product shape (single \texttt{iter\_arg}, two memref loads at the induction variable, multiply--add--yield) and tags it with \texttt{linalg\_op = "dot\_product"}. \\
\addlinespace[1pt]
D8 \texttt{emit-taidl-metadata}
& Walks the lifted module to classify each \texttt{memref} argument by its load/store footprint, label scalar arguments as control attributes, infer grid dimensions from coordinate suffixes in target ASV names, and emit a closed set of \texttt{taidl.*} attributes consumed by Stage~3. \\
\bottomrule
\end{tabular}
\end{table}

\paragraph{Annotation-driven, accelerator-agnostic.}
With the sole exception of loop reconstruction (C6), Phases~B--D do not rewrite the IR but instead attach \texttt{tensorlift.*} attributes to existing operations; the next pass keys off those attributes.
This decoupling lets each pass be verified independently against MLIR's verifier and against the LLVM IR oracle (\S\ref{sec:evaluation}), and lets unrecognized dataflow degrade gracefully: an instruction whose body fails to match the MAC or clamp template simply does not receive the corresponding annotation, and the Stage~3 assembler falls back to an opaque \texttt{copy} or \texttt{unknown} semantics rather than producing incorrect TAIDL.
None of the eight passes contains any Gemmini- or VTA-specific identifier: the structural cues they rely on (sign-extension chains, MAC fan-in, fixed-point clamp idioms, accumulator def-use chains, dot-product loop shape, and RTL signal-name suffixes such as \texttt{pool\_*} or \texttt{im2col\_*}) are properties of common synthesizable RTL patterns rather than of any particular ISA.
The only per-instruction input is the control vector consumed by Pass~B4, which is itself derived automatically (\S\ref{subsec:stage1}); the same \texttt{tensorlift-opt} binary lifts Gemmini's PE, its DMA controllers, and VTA's TensorGemm and TensorAlu without modification.

On the running example, the eight passes together reduce the Gemmini PE update function from 686 lines of bit-level MLIR to 49 lines, of which only 17 encode the actual computation core \texttt{clamp(dot(\%A,\%B)+\%C)}; the remaining lines are arguments, loads, stores, and the \texttt{taidl.*} metadata that Stage~3 turns into a TAIDL specification (\S\ref{sec:evaluation}).

\subsection{Stage 3: TAIDL Assembly}
\label{subsec:stage3}

Stage~3 is a generic Python emitter that turns the artifacts of Stage~2 into a complete TAIDL specification of the form shown in Listing~\ref{lst:taidl-example}.
It consumes two parallel inputs: the \texttt{taidl.*}-annotated MLIR produced by Pass~D8, and a \texttt{state\_machine.json} file produced by an auxiliary script that reads autoGenILA's per-ASV LLVM IR independently of the lifting pipeline.
Both inputs are derived automatically from Stage~1's output; Stage~3 itself contains no per-accelerator code paths.

\paragraph{Assembling the tensor semantics.}
The assembler operates on a closed vocabulary of thirteen \texttt{taidl.*} attribute keys.
After parsing each annotated function into a Python record, it groups buffers across all functions by the triple \emph{(category, type, width)} to derive one \texttt{add\_data\_model} declaration per signature, then merges the per-(instruction, ASV) functions emitted by Stage~1 back into per-instruction groups via a four-level priority chain---an explicit \texttt{--instr-names} list, the original \texttt{instr.txt}, path-based discovery of \texttt{instr.txt} alongside the input MLIR, and finally a longest-prefix heuristic that requires at least three sibling names to commit a prefix---so the same code path handles Gemmini's four-instruction \texttt{ExecuteController} and VTA's single-instruction modules without modification.
Within each group it selects the metadata record with the richest dataflow as a template and dispatches its \texttt{tensor\_op} through a small table that maps each value to its XLA-HLO-style template: \texttt{mac}/\texttt{dot\_product} expand to \texttt{convert+dot+add} (with optional \texttt{clamp}), \texttt{pool} to \texttt{reduce(max)}, \texttt{im2col\_matmul} to \texttt{reshape+dot+add}, and \texttt{copy} to a buffer-to-buffer move.
A second sweep over indexed signal bases (\texttt{strides\_0/1/2}, \texttt{scales\_0/1/2}, $\ldots$) emits \texttt{IF/ELSE} bank-selection bodies for multi-bank configuration instructions, which is what allows TensorLift's specification to capture the three independently configurable DMA banks that the hand-written Gemmini TAIDL collapses into a single \texttt{@a.stride} parameter.
The output is a single Python file that calls the TAIDL API to register data models, instructions, semantics, and (when a state machine is available) initial states, constraints, and updates; adding a new tensor-op family requires extending only the dispatch table and the \texttt{taidl.tensor\_op} vocabulary in Pass~D8.

\paragraph{Inferring the instruction state machine.}
The constraints and updates that gate instructions---e.g., that \texttt{compute\_preloaded} can fire only after \texttt{preload} has armed the controller---are not visible in the Stage~2 output, which is per-instruction by construction.
We recover them with an auxiliary script that reads the same per-ASV update functions in their LLVM IR form, where the \texttt{icmp eq} comparisons against control registers that Pass~A1 collapses are still intact.
The script auto-detects the FSM register by scoring all multi-writer scalar ASVs against three predicates: small bit-width ($\leq\!8$), the presence of \texttt{icmp eq} comparisons \emph{at the register's own bit-width} (which filters out boolean event flags), and at least two distinct compared constants.
Once the FSM register is fixed, each instruction's constraint is the set of constants its update functions compare the register against, with pipeline-flush boundary values filtered as boilerplate; instructions that write the FSM but test no explicit value are handled by a binary-FSM complement rule, in which a writer of the inverse state is gated on the negation of the active state.
Updates are derived symmetrically from non-identity register writes: a write of an \texttt{rs1}/\texttt{rs2} command field into a configuration register becomes \texttt{@s.<reg> = @a.<reg>}, an FSM transition becomes \texttt{@s.<fsm> = <new\_state>}, and any write whose right-hand side is neither a constant nor a single instruction argument is dropped---its semantics already live in the tensor body emitted by the assembler.
The script writes \texttt{state\_machine.json}, which the assembler ingests in its final emission step.

On the running Gemmini ExecuteController example, this two-pronged Stage~3 reproduces the structure of the hand-written reference TAIDL specification: the four extracted instructions (\texttt{config\_ex}, \texttt{preload}, \texttt{compute\_preloaded}, \texttt{compute\_accumulated}) carry the same constraints and the same state-update relations as the reference, with the only difference being that the inferred FSM register is named \texttt{control\_state} (its RTL identifier) where the hand-written specification names it \texttt{preload} (a TAIDL source-level alias).
This is purely a name choice in the generated TAIDL source: as we report in \S\ref{sec:evaluation}, the test oracles produced from TensorLift's specification are character-identical to those produced from the hand-written reference, because the oracle generator alpha-renames state registers consistently across both inputs.
Together with Stages~1--2, the pipeline closes the loop from RTL to ACT-consumable TAIDL, replacing what is otherwise a hand-authoring step with a deterministic lifting and assembly process.
